% AUTO-GENERATED by the update_record tool. DO NOT EDIT.
% verify_paper regenerates this file from record.json and the run outputs
% and fails on any difference, so manual edits always surface.
\noindent\textbf{H1.} 画像分類において、ラベルスムージング（α=0.1）で学習した教師から低温度（T=1）で蒸留した生徒は、hard target で学習した教師から同条件で蒸留した生徒より較正が良い（期待較正誤差 ECE が低い）。
\emph{Grounded on:} \texttt{\detokenize{s1.p1}}, \texttt{\detokenize{s1.p2}}, \texttt{\detokenize{s3.p2}}, \texttt{\detokenize{s4.p1}}, \texttt{\detokenize{s4.p2}}.
\begin{enumerate}
\item[\textbf{C1}] LS 教師から T=1 で蒸留した生徒の ECE は、hard 教師から T=1 で蒸留した生徒の ECE より 0.01 以上低い。
  \emph{Rationale:} 仮説そのものの測定である。教師の較正差（s1.p1）が蒸留を通じて生徒に現れるか（s4.p1）を、Chandrasegaran らが精度の面で推奨する T=1（s3.p2）で直接比較する。
  \emph{Cites:} \texttt{\detokenize{s1.p1}}, \texttt{\detokenize{s3.p2}}, \texttt{\detokenize{s4.p1}}; \texttt{\detokenize{d1}}: \texttt{\detokenize{s3.p2}}; \texttt{\detokenize{student-ls}}: \texttt{\detokenize{s3.p2}}.
  \emph{Criterion:} \texttt{\detokenize{student-ls}}.\texttt{\detokenize{ece}} $-$ \texttt{\detokenize{student-hard}}.\texttt{\detokenize{ece}} $\leq -0.01$.
  \emph{Prediction:} $[-0.04, -0.01]$ (Müller らは CIFAR-100 で LS が教師の ECE を数ポイント下げると報告しており、その一部が生徒に継承されるという想定。).
  \emph{Observed:} pending.
  \emph{Verdict:} pending.
\end{enumerate}
\noindent\emph{Assumed, so that the claims together imply H1:}
\begin{itemize}
\item 15-bin ECE が「較正の良さ」を代表する（c1）。
\item CIFAR-100 と ResNet-18 教師 / 小型 CNN 生徒という 1 設定での差が、画像分類一般に外挿できる（c1）。
\item hard target 教師からの蒸留が「較正を継承しない」側の代表的な基準になっている（c1）。
\item seed 1 本の ECE 差 0.01 が run 間のばらつきより大きい（c1）。
\end{itemize}
\noindent\textbf{Sources.}
\noindent\textbf{S1} \texttt{\detokenize{muller-2019-when}}: When Does Label Smoothing Help? (2019).
\begin{itemize}
\item[\texttt{\detokenize{s1.p1}}] (result) ``label smoothing improves model calibration which can significantly improve beam-search''
\item[\texttt{\detokenize{s1.p2}}] (result) ``if a teacher network is trained with label smoothing, knowledge distillation into a student network is much less effective''
\end{itemize}
\noindent\textbf{S2} \texttt{\detokenize{shen-2021-label}}: Is Label Smoothing Truly Incompatible with Knowledge Distillation: An Empirical Study (2021).
\begin{itemize}
\item[\texttt{\detokenize{s2.p1}}] (claim) ``label smoothing erases relative information between teacher logits''
\end{itemize}
\noindent\textbf{S3} \texttt{\detokenize{chandrasegaran-2022-revisiting}}: Revisiting Label Smoothing and Knowledge Distillation Compatibility: What was Missing? (2022).
\begin{itemize}
\item[\texttt{\detokenize{s3.p1}}] (result) ``This systematic diffusion essentially curtails the benefits (as claimed by Shen et al. (2021b)) obtained by distilling from an LS-trained teacher, thereby rendering KD at increased temperatures ineffective.''
\item[\texttt{\detokenize{s3.p2}}] (method) ``We suggest to use an LS-trained teacher with a low-temperature transfer (i.e. T = 1) to achieve high performance students.''
\end{itemize}
\noindent\textbf{S4} \texttt{\detokenize{sultan-2023-knowledge}}: Knowledge Distillation ≈ Label Smoothing: Fact or Fallacy? (2023).
\begin{itemize}
\item[\texttt{\detokenize{s4.p1}}] (result) ``In KD, the student inherits not only its knowledge but also its confidence from the teacher''
\item[\texttt{\detokenize{s4.p2}}] (setup) ``Experiments on four text classification tasks involving models of different sizes''
\end{itemize}
